%----------------------------------------------------------------------------------------
%    PACKAGES
%----------------------------------------------------------------------------------------

\documentclass[aspectratio=169,xcolor=dvipsnames]{beamer}
\usetheme{SimpleDarkBlue}
\usepackage{hyperref}
\usepackage{graphicx} % Allows including images
\usepackage{booktabs} % Allows the use of \toprule, \midrule and \bottomrule in tables

%----------------------------------------------------------------------------------------
%    TÍTULO
%----------------------------------------------------------------------------------------

\title{Análisis de la relación entre el uso de redes sociales y el rendimiento académico en estudiantes}

\author{Sebastián Calderón, Luis Diego Elizondo, Jeremy Garcia, Katia Moreno}

\date{\today}

\begin{document}    

%----------------------------------------------------------------------------------------
%    PRESENTACIÓN
%----------------------------------------------------------------------------------------

\begin{frame}[plain,noframenumbering]
    \titlepage
    \begin{figure}
    
\includegraphics[width=0.2\linewidth]{logoucr.png}
    \end{figure}
\end{frame}

%-----------------------------------------------------------------------------------------

\begin{frame}{Pregunta del proyecto}
\centering

\Large
\textbf{¿Existe relación entre el uso de redes sociales y el rendimiento académico?}


\vspace{1cm}
    \begin{alertblock}{Problemas}
    \begin{itemize}
        \item Uso de redes sociales con fines académicos.
        \item Posibles problemas personales no considerados.
    \end{itemize}
        
    \end{alertblock}

    
\end{frame}

%-----------------------------------------------------------------------------------------

\begin{frame}{Datos}
\begin{block}{Tipo de las variables}
\centering
\begin{tabular}{c c}
\toprule
\textbf{Variable} & \textbf{Tipo} \\
\midrule
Promedio académico (GPA) & Numérica \\ 
Horas diarias en redes sociales & Numérica \\ 
Horas diarias de estudio & Numérica \\ 
Porcentaje de asistencia & Numérica \\ 
Edad & Numérica \\ Distracción & Numérica \\
Género & Categórica \\ 
Red social principal & Categórica \\ 
Uso nocturno de redes & Categórica \\ 
Zona de residencia (urbano/rural) & Categórica \\ 
\bottomrule
\end{tabular}
\end{block}
\vspace{2mm}
\centering
{\small Variables principales utilizadas en el análisis}

\end{frame}

%-----------------------------------------------------------------------------------------

\begin{frame}{Análisis exploratorio de datos}
  \begin{columns}

    \begin{column}{0.45\textwidth}
      \centering
      \includegraphics[width=\linewidth]{grafico_1_horas_redes_vs_promedio.png}
      \captionof{(a) Horas en redes sociales vs Promedio académico} % Opcional
    \end{column}
    
    \hspace{0.05\textwidth}
    
    \begin{column}{0.45\textwidth}
      \centering
      \includegraphics[width=\linewidth]{grafico_3_horas_estudio_vs_promedio.png}
      \captionof{(b) Horas de estudio vs Promedio académico} 
    \end{column}
  \end{columns}
\end{frame}

%-----------------------------------------------------------------------------------------

\begin{frame}{Análisis exploratorio de datos}
  \begin{columns}

    \begin{column}{0.45\textwidth}
      \centering
      \includegraphics[width=\linewidth]{grafico_2_categoria_redes_vs_promedio.png}
      \captionof{(c) Distribución del promedio por nivel de uso de redes} % Opcional
    \end{column}
    
    \hspace{0.05\textwidth}
    
    \begin{column}{0.45\textwidth}
      \centering
      \includegraphics[width=\linewidth]{grafico_4_promedio_por_plataforma.png}
      \captionof{(d) Promedio académico medido por red social principal} 
    \end{column}
  \end{columns}
\end{frame}

%-----------------------------------------------------------------------------------------


\begin{frame}{Análisis exploratorio de datos}

\begin{center}
    \includegraphics[width=0.58\textwidth]{grafico_5_heatmap_redes_distraccion_promedio.png}

    \vspace{0.25cm}

    \textbf{(e) Uso de redes y distracción académica}
\end{center}

\end{frame}

%-----------------------------------------------------------------------------------------

\begin{frame}{Metodología propuesta}

\vspace{0.25cm}

\begin{columns}[T]

\column{0.48\textwidth}

\textbf{Análisis comparativo}

\begin{itemize}
    \item Prueba t para comparar promedio entre bajo uso y alto uso de redes.
    \item ANOVA para comparar promedio entre varias categorías de uso.
\end{itemize}

\vspace{0.2cm}


\column{0.48\textwidth}

\textbf{Análisis de asociación}

\begin{itemize}
    \item Correlación de Pearson o Spearman entre horas en redes y promedio.
\end{itemize}




\end{columns}

\vspace{0.25cm}

\begin{alertblock}{Alternativa descartada}
Consideramos clasificar por red social principal, especialmente YouTube por su posible uso académico,
pero lo descartamos como eje central porque la pregunta principal se responde mejor desde la intensidad de uso,
no desde la plataforma.
\end{alertblock}

\end{frame}

%-----------------------------------------------------------------------------------------

%-----------------------------------------------------------------------------------------

\begin{frame}{Resultados esperados}

\small

\begin{columns}[T]

\column{0.52\textwidth}

\textbf{Patrones esperados}

\begin{itemize}
    \item Diferencias en rendimiento académico entre grupos con distinto uso de redes sociales.
    \item Mayor dispersión y promedios ligeramente menores en estudiantes con uso elevado.
    \item Asociación positiva entre horas de estudio y rendimiento académico.
\end{itemize}

\column{0.42\textwidth}

\begin{block}{Idea central}
No esperamos necesariamente una relación lineal fuerte, sino una asociación condicionada
por la distracción académica y las horas de estudio.
\end{block}

\end{columns}

\vspace{0.3cm}

\begin{alertblock}{Criterio para responder la pregunta}
Consideraremos respondida la pregunta si las pruebas estadísticas y el modelo permiten determinar
si el uso de redes sociales, solo o combinado con distracción académica, se asocia significativamente
con el promedio académico.
\end{alertblock}

\end{frame}

%-----------------------------------------------------------------------------------------

\end{document}
